In our runtime, every value starts with a header containing two
tags. The first tag specifies the value kind: \kw{ctor}, \kw{pap},
\kw{array}, \kw{string}, \kw{num}, \kw{thunk}, or \kw{task}. The
second tag specifies whether the value is single-threaded,
multi-threaded, or \emph{persistent}. We will describe how this kind is used to
implement thread safe reference counting in the next subsection. The kinds \kw{ctor} and
\kw{pap} are used to implement the corresponding values used in the
formal semantics of \pureIR{} and \rcIR{}.
The kinds \kw{array}, \kw{string}, and \kw{thunk} are self
explanatory.  The kind \kw{num} is for arbitrary precision numbers
implemented using the GNU multiple precision library (GMP).
The \kw{task} value is described in the next subsection.

Values tagged as single- or multi-threaded also contain a reference
counter. This counter is stored in front of the standard value header.
We will primarily focus on the layout of \kw{ctor} values here because  it is the most
relevant one for the ideas presented in this paper. A $\kw{ctor}_i$
value header also includes the constructor index $i$, the number of
pointers to other values and/or boxed values, and the number of bytes used to store scalar
unboxed values such as machine integers and enumeration types. In a
64-bit machine, the \kw{ctor} value header is 16 bytes long, twice the
size of the header used in OCaml to implement the corresponding kind
of value.  After the header, we store all pointers to other values
and boxed values, and then all unboxed values. Thus, in a 64 bit machine,
our runtime uses 32 bytes to implement a \emph{List} \emph{Cons} value:
16 bytes for the header, and 16 bytes for storing the list head and tail.
The unboxed value support has restrictions similar to the ones found in GHC.
For example, to pass an unboxed value to a polymorphic function we must first box it.

Our runtime has built-in support for array and string operations.
Strings are just a special case of arrays where the elements are
characters. We perform destructive updates when the array is not shared.
For example, given the array write primitive
\begin{lstlisting}
Array.write : Array $\alpha$ -> Nat -> $\alpha$ -> Array $\alpha$
\end{lstlisting}
the function application $\textit{Array.write}\ a\ i\ v$ will
destructively update and return the array $a$ if it is not shared.
This is a well known optimization for systems based on
reference counting~\cite{GCbook}, nonetheless we mention it here
because it is relevant for many applications. Moreover,
destructive array updates and our \kw{reset}/\kw{reuse} technique
complement each other. As an example, if we have a nonshared list of
integer arrays \textit{xs},
$\textit{map}\ (\textit{Array.map}\ \textit{inc})\ \textit{xs}$
destructively updates the list and all arrays. In the experimental
section we demonstrate that our pure quick sort is as efficient as the
quick sort using destructive updates in OCaml, and the quick sort
using the primitive ST monad in Haskell.

%% in automated reasoning and proof automation,
%% where hash-tables and arrays are frequently used in the implementation
%% of decision procedures.

%% When we create array and string values, we allocate extra storage to make sure
%% primitives such as
%% \begin{align*}
%%   \textit{array.push} &: \textit{Array}\ \alpha \rightarrow \alpha \rightarrow \textit{Array}\ \alpha \\
%%   \textit{string.push} &: \textit{String} \rightarrow \textit{Char} \rightarrow \textit{String} \\
%%   \textit{string.append} &: \textit{String} \rightarrow \textit{String} \rightarrow \textit{String}
%% \end{align*}
%% do not necessarily reallocate memory each time they are used when the
%% reference counter of the input value is 1. This approach is similar
%% to the \emph{vector} class used in imperative languages such as C++
%% and Java.

%%% Local Variables:
%%% mode: latex
%%% TeX-master: "refcount"
%%% End:
